\chapter{The Jacobian of a smooth proper curve}
\label{chap:Challenge}

Throughout, \(k\) is a field and \(C\) a \emph{smooth curve} over \(k\): a geometrically
irreducible, smooth \(k\)-scheme of relative dimension one, proper over \(k\). We write
\(C \to \Spec k\) for its structure morphism. This chapter states the goal of the project — the
genus, the Jacobian, the Abel--Jacobi map, and the universal property that characterizes the
Jacobian as the Albanese variety of \(C\). It is the mathematical counterpart of the statement file
\texttt{AlgebraicJacobian/Challenge.lean}.

\emph{Provenance.} Each node below should carry a \verb!\source{...}! anchor to the reference it is
read from once that reference has been consulted; anchors are added during the rebuild, never
guessed.

\section{Genus}

\begin{definition}[Genus]
  \label{def:genus}
  \lean{AlgebraicGeometry.genus}
  \leanok
  \uses{def:moduleKSheaf, def:HModule}
  The \emph{genus} of a smooth proper curve \(C\) over \(k\) is
  \[
    g(C) \;=\; \dim_k H^1\!\bigl(C, \struct C\bigr),
  \]
  the \(k\)-dimension of the first cohomology group of the structure sheaf, viewed as a sheaf
  of \(k\)-modules on the small Zariski site of \(C\) (Definition~\ref{def:moduleKSheaf}),
  with cohomology the \(\Ext\)-groups from the constant sheaf
  (Definition~\ref{def:HModule}).
\end{definition}

\section{The Jacobian as an abelian variety}

\begin{definition}[Jacobian]
  \label{def:jacobian}
  \lean{AlgebraicGeometry.Jacobian}
  The \emph{Jacobian} \(\Jac(C)\) of a smooth proper curve \(C\) over \(k\) is the identity
  component \(\Pic^0_{C/k}\) of the relative Picard scheme of \(C\) over \(k\): the \(k\)-scheme
  representing (the fppf-sheafification of) the functor of degree-zero line bundles on \(C\).
\end{definition}

\begin{proposition}[Group scheme structure]
  \label{prop:jacobian_grpobj}
  \uses{def:jacobian}
  \lean{AlgebraicGeometry.Jacobian.instGrpObj}
  \(\Jac(C)\) is a commutative group scheme over \(k\), with multiplication induced by the tensor
  product of line bundles.
\end{proposition}

\begin{theorem}[Smoothness of the Jacobian]
  \label{thm:jacobian_smooth}
  \uses{def:genus, def:jacobian}
  \lean{AlgebraicGeometry.Jacobian.smoothOfRelativeDimension_genus}
  \(\Jac(C)\) is smooth over \(k\) of relative dimension \(g(C)\).
\end{theorem}

\begin{theorem}[Properness of the Jacobian]
  \label{thm:jacobian_proper}
  \uses{def:jacobian}
  \lean{AlgebraicGeometry.Jacobian.instIsProper}
  \(\Jac(C)\) is proper over \(k\).
\end{theorem}

\begin{theorem}[Geometric irreducibility of the Jacobian]
  \label{thm:jacobian_geomirred}
  \uses{def:jacobian}
  \lean{AlgebraicGeometry.Jacobian.instGeometricallyIrreducible}
  \(\Jac(C)\) is geometrically irreducible over \(k\).
\end{theorem}

In particular \(\Jac(C)\) is an abelian variety over \(k\) of dimension \(g(C)\).

\section{The Abel--Jacobi map and the universal property}

\begin{definition}[Abel--Jacobi map]
  \label{def:ofCurve}
  \uses{def:jacobian}
  \lean{AlgebraicGeometry.Jacobian.ofCurve}
  For a \(k\)-rational point \(P : \Spec k \to C\), the \emph{Abel--Jacobi map}
  \(\varphi_P : C \to \Jac(C)\) sends a point \(Q\) to the class of the line bundle
  \(\struct C(Q - P)\).
\end{definition}

\begin{lemma}[Base point maps to the origin]
  \label{lem:comp_ofCurve}
  \uses{def:ofCurve, prop:jacobian_grpobj}
  \lean{AlgebraicGeometry.Jacobian.comp_ofCurve}
  The Abel--Jacobi map \(\varphi_P\) sends the base point \(P\) to the neutral element of
  \(\Jac(C)\): \(\varphi_P \circ P = 0\).
\end{lemma}

\begin{theorem}[Universal property of the Jacobian]
  \label{thm:exists_unique_ofCurve_comp}
  \uses{def:ofCurve, lem:comp_ofCurve, prop:jacobian_grpobj}
  \lean{AlgebraicGeometry.Jacobian.exists_unique_ofCurve_comp}
  \(\Jac(C)\) is the Albanese variety of \(C\): for every abelian variety \(A\) over \(k\) and every
  morphism \(f : C \to A\) with \(f \circ P = 0\), there is a unique homomorphism of group schemes
  \(g : \Jac(C) \to A\) such that \(f = g \circ \varphi_P\).
\end{theorem}
